% The instrument in the supply path, with three markers on one timebase.
% Coordinates are explicit so that no edge crosses a block.
\begin{tikzpicture}[font=\sffamily\small]
% ---------------- the cycle, as a column
\node[fw, text width=30mm] (wake)  at (0,0)     {wake\\{\scriptsize HSI, PLL1 off}};
\node[fw, text width=30mm] (rest)  at (0,-1.35) {restore\\{\scriptsize scale, latency, PLL}};
\node[fw, text width=30mm] (sense) at (0,-2.7)  {sense\\{\scriptsize one sample, one feature}};
\node[fw, text width=30mm] (prep)  at (0,-4.05) {prepare\\{\scriptsize markers low, tick off}};
\node[hw, text width=30mm] (stop)  at (0,-5.4)  {stop\\{\scriptsize SLEEPDEEP, WFI}};
\begin{scope}[on background layer]
\node[layer, fit=(wake)(rest)(sense)(prep)(stop), inner xsep=7pt, inner ysep=10pt,
      label={[layerlabel, anchor=south west]north west:one cycle}] (brd) {};
\end{scope}
\draw[flow] (wake) -- (rest);
\draw[flow] (rest) -- (sense);
\draw[flow] (sense) -- (prep);
\draw[flow] (prep) -- (stop);

% ---------------- what keeps time while the bus clocks are stopped
\node[hw, text width=22mm] (lse)   at (-6.2,-1.0) {32.768 kHz crystal\\{\scriptsize fitted and confirmed}};
\node[hw, text width=22mm] (rtc)   at (-6.2,-2.6) {RTC alarm};
\node[hw, text width=22mm] (lptim) at (-6.2,-4.2) {LPTIM1};
\begin{scope}[on background layer]
\node[layer, fit=(lse)(rtc)(lptim), inner xsep=7pt, inner ysep=10pt,
      label={[layerlabel, anchor=south west]north west:still running}] {};
\end{scope}
\draw[flow] (lse) -- (rtc);
\draw[flow] (lse.west) -- (-7.8,-1.0) -- (-7.8,-4.2) -- (lptim.west);

% ---------------- the wake line, gathered onto one spine
\draw[draw=red!70!black, dashed, line width=0.8pt]
   (rtc.east) -- (-4.3,-2.6) -- (-4.3,-4.2) -- (lptim.east);
\draw[irq] (-4.3,-2.6) -- (-4.3,1.5)
   -- node[lbl, above]{wake line} (0,1.5) -- (wake.north);

% ---------------- the alarm that closes the cycle
\draw[flow] (stop.west) -- (-2.6,-5.4)
   -- node[lbl, rotate=90]{about 10 s} (-2.6,0) -- (wake.west);

% ---------------- the instrument and the host
\node[ext, text width=32mm] (ppk)  at (5.9,-2.7) {nRF-PPK2\\{\scriptsize ammeter, 1 A maximum}};
\node[ext, text width=32mm] (host) at (5.9,-4.6) {host\\{\scriptsize capture.py, cycles.py}};
\begin{scope}[on background layer]
\node[layer, fit=(ppk)(host), inner xsep=7pt, inner ysep=10pt,
      label={[layerlabel, anchor=south west]north west:the instrument}] {};
\end{scope}
\draw[bus] (brd.east) -- node[lbl, above]{IDD and markers} (ppk.west);
\draw[flow] (ppk) -- node[lbl, right]{4 channels} (host);

\node[umlnote, text width=52mm, anchor=north west] at (-8.0,-6.4)
  {The probe is still powered from USB and is not in the measured path. Which jumper carries this current and what it isolates are items 6 and 34 on the confirm list.};
\node[umlnote, text width=52mm, anchor=north west] at (1.9,-6.4)
  {Current and three marker pins arrive on one timebase, so the trace can be segmented without a human looking at it.};
\end{tikzpicture}
